\section{USAMO 1984}
        \begin{enumerate}
        	\item (\textit{USAMO 1984}) The product of two of the four roots of the quartic equation $x^4 - 18x^3 + kx^2+200x-1984=0$ is $-32$. Determine the value of $k$.


 



\textbf{Solution 1 (ingenious)}

Using Vieta's formulas, we have: 

 
 \begin{align*}a+b+c+d &= 18,\\ ab+ac+ad+bc+bd+cd &= k,\\ abc+abd+acd+bcd &=-200,\\ abcd &=-1984.\\ \end{align*}
 

 
 From the last of these equations, we see that $cd = \frac{abcd}{ab} = \frac{-1984}{-32} = 62$. Thus, the second equation becomes $-32+ac+ad+bc+bd+62=k$, and so $ac+ad+bc+bd=k-30$. The key insight is now to factor the left-hand side as a product of two binomials: $(a+b)(c+d)=k-30$, so that we now only need to determine $a+b$ and $c+d$ rather than all four of $a,b,c,d$. 

 Let $p=a+b$ and $q=c+d$. Plugging our known values for $ab$ and $cd$ into the third Vieta equation, $-200 = abc+abd + acd + bcd = ab(c+d) + cd(a+b)$, we have $-200 = -32(c+d) + 62(a+b) = 62p-32q$. Moreover, the first Vieta equation, $a+b+c+d=18$, gives $p+q=18$. Thus we have two linear equations in $p$ and $q$, which we solve to obtain $p=4$ and $q=14$. 

 Therefore, we have $(\underbrace{a+b}_4)(\underbrace{c+d}_{14}) = k-30$, yielding $k=4\cdot 14+30 = \boxed{86}$.

 



\textbf{Solution 2 (cool)}

We start as before: $ab=-32$ and $cd=62$. We now observe that a and b must be the roots of a quadratic, $x^2+rx-32$, where r is a constant (secretly, r is just -(a+b)=-p from Solution \#1). Similarly, c and d must be the roots of a quadratic $x^2+sx+62$. 

 Now 

 
 \begin{align*}x^4-18x^3+kx^2+200x-1984 =& (x^2+rx-32)(x^2+sx+62)\\  =& x^4+(r+s)x^3+(62-32+rs)x^2\\ &+(62s-32r)x-1984.\end{align*}


 Equating the coefficients of $x^3$ and $x$ with their known values, we are left with essentially the same linear equations as in Solution \#1, which we solve in the same way. Then we compute the coefficient of $x^2$ and get $k=\boxed{86}.$

 



\textbf{Solution 3 AM-GM}

Let $r_1, r_2, r_3, r_4$ be the roots of the polynomial $x^4 - 18x^3 + kx^2 + 200x - 1984 = 0$. We are given that $r_1 r_2 = -32$.

 By Vieta's formulas, we have:

 $r_1+r_2+r_3+r_4 = 18$

 $r_1r_2+r_1r_3+r_1r_4+r_2r_3+r_2r_4+r_3r_4 = k$

 $r_1r_2r_3+r_1r_2r_4+r_1r_3r_4+r_2r_3r_4 = -200$

 $r_1r_2r_3r_4 = -1984$

 Since $r_1r_2 = -32$, we have $(-32)r_3r_4 = -1984$, so $r_3r_4 = \frac{-1984}{-32} = 62$.

 Also, $r_1r_2r_3+r_1r_2r_4+r_1r_3r_4+r_2r_3r_4 = -200$, so $r_1r_2(r_3+r_4)+r_3r_4(r_1+r_2) = -200$. Substituting $r_1r_2 = -32$ and $r_3r_4 = 62$, we have $-32(r_3+r_4)+62(r_1+r_2) = -200$.

 Let $A = r_1+r_2$ and $B = r_3+r_4$. Then $A+B = 18$, so $B = 18-A$. Substituting this into the equation, we have $-32B+62A = -200$, so $-32(18-A)+62A = -200$. $-576+32A+62A = -200$, so $94A = 376$, which means $A = \frac{376}{94} = 4$. Then $B = 18-A = 18-4 = 14$.

 So we have $r_1+r_2 = 4$ and $r_1r_2 = -32$. Then $r_1$ and $r_2$ are roots of $x^2-4x-32 = 0$, so $(x-8)(x+4) = 0$, which means $r_1 = 8$ and $r_2 = -4$ (or vice versa).

 Also we have $r_3+r_4 = 14$ and $r_3r_4 = 62$. Then $r_3$ and $r_4$ are roots of $x^2-14x+62 = 0$. Using the quadratic formula, $x = \frac{14 \pm \sqrt{14^2-4(62)}}{2} = \frac{14 \pm \sqrt{196-248}}{2} = \frac{14 \pm \sqrt{-52}}{2} = 7 \pm i\sqrt{13}$. Then $r_3 = 7+i\sqrt{13}$ and $r_4 = 7-i\sqrt{13}$.

 We want to find $k = r_1r_2+r_1r_3+r_1r_4+r_2r_3+r_2r_4+r_3r_4$. $k = r_1r_2+r_3r_4 + (r_1+r_2)(r_3+r_4) = -32+62+(4)(14) = 30+56 = 86$.

 Final Answer: The final answer is $\boxed{86}$

 



\textbf{Solution 4 (Alcumus)}

Since two of the roots have product $-32,$ the equation can be factored in the form
 \[x^4 - 18x^3 + kx^2 + 200x - 1984 = (x^2 + ax - 32)(x^2 + bx + c).\]Expanding, we get
 \[x^4 - 18x^3 + kx^2 + 200x - 1984 = x^4 + (a + b) x^3 + (ab + c - 32) x^2 + (ac - 32b) x - 32c = 0.\]Matching coefficients, we get\begin{align*}a + b &= -18, \\ab + c - 32 &= k, \\ac - 32b &= 200, \\-32c &= -1984.\end{align*}
Then $c = \frac{-1984}{-32} = 62,$ so $62a - 32b = 200.$ With $a + b = -18,$ we can solve to find $a = -4$ and $b = -14.$ Then
 \[k = ab + c - 32 = \boxed{86}.\]

 



\textbf{Solution 5 (If You Don't Know Vieta)}

We know that any given quartic polynomial can be factored into two quadratic polynomials. We see that for any given quadratic, the product of it's roots is:
 \[\frac{-b - \sqrt{b^{2}-4ac}}{2a} \cdot \frac{-b + \sqrt{b^{2}-4ac}}{2a} = \frac{b^{2}-b^{2}+4ac}{4a^{2}} = \frac{c}{a}\]We see that both quadratics must have leading coefficient of $1$, so $a = 1$

 We first try a quadratic where the product of it's roots is $-32$
 \[x^{2}+b_{1}x-32\]We then see that the other quadratic has the form
 \[x^{2}+b_{2}x+62\]since the constant term in both quadratics must multiply to $-1984$

 To get our general form quartic from these, we multiply
 \[(x^{2}+b_{1}x-32)(x^{2}+b_{2}x+62)\]Expanding, then combining coefficients, we get
 \[x^{4}+b_{2}x^{3}+62x^{2}+b_{1}x^{3}+b_{1}b_{2}x^{2}+62b_{1}x-32x^{2}-32b_{2}x-1984\]
 \[x^{4}+(b_{1}+b_{2})x^{3}+(30+b_{1}b_{2})x^{2}+(62b_{1}-32b_{2})x-1984\]Looking back at our original polynomial, we extract the system
 \[\begin{cases} b_{1}+b_{2} = -18 \\ 62b_{1}-32b_{2} =200 \\ 30+b_{1}b_{2} = k \end{cases}\]To solve it, we see
 \[b_{1} =-b_{2}-18\]
 \[62(-b_{2}-18)-32b_{2} =200\]
 \[31(-b_{2}-18)-16b_{2} =100\]
 \[-47b_{2} =658\]
 \[b_{2} = -14\]
 \[b_{1} = -4\]
 \[k = 30 + 56\]
 \[\boxed{k=86}\]~shockfront99

 


	\item (\textit{USAMO 1984}) The geometric mean of any set of $m$ non-negative numbers is the $m$-th root of their product.


 $\quad (\text{i})\quad$ For which positive integers $n$ is there a finite set $S_n$ of $n$ distinct positive integers such that the geometric mean of any subset of $S_n$ is an integer?


 $\quad (\text{ii})\quad$ Is there an infinite set $S$ of distinct positive integers such that the geometric mean of any finite subset of $S$ is an integer?


 



\textbf{Solution}

a) We claim that for any numbers $p_1$, $p_2$, ... $p_n$, $p_1^{n!}, p_2^{n!}, ... p_n^{n!}$ will satisfy the condition, which holds for any number $n$.

 Since $\sqrt[n] ab = \sqrt[n] a * \sqrt[n] b$, we can separate each geometric mean into the product of parts, where each part is the $k$th root of each member of the subset and the subset has $k$ members.

 Assume our subset has $k$ members. Then, we know that the $k$th root of each of these members is an integer (namely $p^{n!/k}$), because $k \leq n$ and thus $k  |  n!$. Since each root is an integer, the geometric mean will also be an integer.

 b) If we define $q$ as an arbitrarily large number, and $x$ and $y$ as numbers in set $S$, we know that ${\sqrt[q]{\frac{x}{y}}}$ is irrational for large enough $q$, meaning that it cannot be expressed as the fraction of two integers. However, both the geometric mean of the set of $x$ and $q-1$ other arbitrary numbers in $S$ and the set of $y$ and the same other $q-1$ numbers are integers, so since the other numbers cancel out, the geometric means divided, or ${\sqrt[q]{\frac{x}{y}}}$, must be rational. This is a contradiction, so no such infinite $S$ is possible.

 -aops111 (first solution dont bully me)

 



\textbf{Solution 2}

(i) The solution is the same as in the first solution

 (ii)Let $p$ and $q$ be some prime numbers, and let $v_p(n)$ be the largest number such that $p^{v_p(n)}$ divides $n$. $v_p(n)$ is also the exponent of $p$ in the prime factorization of $n$ (or $0$, if $p$ isn't a divisor of $n$). Then, because $S$ is infinite, we can choose $q-1$ numbers $a_1, a_2, ... a_{q-1}$ in $S$ such that $v_p(a_i)$ all give the same remainder when divided by $q$. (Otherwise, because $v_p(n) \mod q$ has $q$ distinct values, $S$ could contain at most $q(p-2)$ numbers, which is finite.) Call that remainder $r$.

 Now take some integer $b$ in $S$, not part of the other $q-1$ values. Together, their product must be a $q$th power, meaning $q|v_p(a_1a_2...a_{q-1}b)$ Thus, $v_p(a_1a_2...a_{q-1}b) \equiv 0 \pmod q$. It can be seen easily that for two positive integers $x$ and $y$ we have $v_p(ab) = v_p(a) + v_p(b)$, and by definition $v_p(a_i) \equiv r \pmod q$, so we can write $v_p(b) + (q-1)r \equiv 0 \pmod q$. Adding $r$ to both sides we have $v_p(b) \equiv v_p(b) + qr \equiv r \pmod q$.

 The choice of $b$ was arbitrary, meaning for all numbers $n$ in $S$, $v_p(n)$ modulo $q$ is the constant $r$, whether part of $a_1, a_2... a_{q-1}$ or not. The choice of $q$ was also arbitrary, thus for two integers $m$ and $n$ in $S$, $v_p(m) = v_p(n)$. Otherwise, we could pick $q$ larger than both values, and get a contradiction. Thus for all  numbers $n$ in $S$, $v_p(n)$ is constant. The choice of $p$ was also arbitrary, meaning any two distinct numbers in $S$ have the same prime factorization. By the fundamental theorem of arithmetic, this means the two numbers are the same, contradicting with them being distinct. Thus, no such $S$ exists.

 -Circling

 


	\item (\textit{USAMO 1984}) $P$, $A$, $B$, $C$, and $D$ are five distinct points in space such that $\angle APB = \angle BPC = \angle CPD = \angle DPA = \theta$, where $\theta$ is a given acute angle. Determine the greatest and least values of $\angle APC + \angle BPD$.


 



\textbf{Solution}

Greatest value is achieved when all the points are as close as possible to all being on a plane.

 Since $\theta < \frac{\pi}{2}$, then $\angle APC + \angle BPD < \pi$

 Smallest value is achieved when point P is above and the remaining points are as close as possible to colinear when $\theta > 0$, then $\angle APC + \angle BPD > 0$

 and the inequality for this problem is:

 $0 < \angle APC + \angle BPD < \pi$

 <i>Alternate solutions are always welcome. If you have a different, elegant solution to this problem, please \href{https://artofproblemsolving.com/wiki/index.php?title=1984\_USAMO\_Problems/Problem\_3&action=edit}{add it} to this page.</i>

 


	\item (\textit{USAMO 1984}) A difficult mathematical competition consisted of a Part I and a Part II with a combined total of $28$ problems. Each contestant solved $7$ problems altogether. For each pair of problems, there were exactly two contestants who solved both of them. Prove that there was a contestant who, in Part I, solved either no problems or at least four problems.


 



\textbf{Solution}

First, notice that for any given problem, among the participants who solved that problem, exactly two of them must have also solved each of the other 27 problems. Since each person solved exactly 6 of the other problems, the number of participants who solved any given problem must be $\frac{2\cdot 27}{6} = 9$.This also allows us to easily calculate the number of participants as $\frac{9\cdot 28}{7} = 36$.

 Now, let $n$ be the number of problems in part I, and assume that no participant solved 0 or more than 4 problems in part I.In other words, assume that all participants solved either 1, 2 or 3 problems in part I. Let $m_1, m_2, m_3$ be the number of participants who solved exactly 1, 2 and 3 problems in part I, respectively.

 Clearly, this implies $m_1+m_2+m_3 = 36$.

 Also, since we know 9 participants solved each problem, $m_1+2m_2+3m_3 = 9n$, as both sides must count the total number of solves in part I.

 Clearly $n$ is at least 4, as 36 participants all solved at least 1 of the problems in part I, but only 9 could have solved each of them.Now consider that exactly two participants must have solved each of the $\binom{n}{2}$ pairs of problems within part I. Since each participant who solved 2 problems in part I solved 1 pair of part I problems, and each participant who solved 3 problems in part I solved 3 pairs of part I problems, we can write $m_2 + 3m_3 = 2\binom{n}{2}$.

 We can now simplify this system of three equations to solve for $m_3, m_2, m_1$ in terms of $n$. If we first solve for $m_3$ we get $m_3 = 2\binom{n}{2} + 36 - 9n = n^2 - 10n + 36$. If we then solve for $m_2$ we get $9n-36 = m_2 + 2m_3 = m_2 + 2(n^2-10n+36)$ and so $m_2 = 29n - 2n^2 - 108$, but the right hand side is negative for all values of $n$, and clearly $m_2$ cannot be negative. Thus we have reached a contradiction.

 


 


	\item (\textit{USAMO 1984}) $P(x)$ is a polynomial of degree $3n$ such that


 
 \begin{eqnarray*} P(0) = P(3) = \cdots &=& P(3n) = 2, \\ P(1) = P(4) = \cdots &=& P(3n-2) = 1, \\ P(2) = P(5) = \cdots &=& P(3n-1) = 0, \quad\text{ and }\\ && P(3n+1) = 730.\end{eqnarray*}


 Determine $n$.


 



\textbf{Solution}

By Lagrange Interpolation Formula $f(x) = 2\sum_{p=0}^{n}\left ( \prod_{0\leq r\neq3p\leq 3n}^{{}}\frac{x-r}{3p-r} \right )+ \sum_{p=1}^{n}\left ( \prod_{0\leq r\neq3p-2\leq 3n}^{{}} \frac{x-r}{3p-2-r}\right )$

 and hence $f(3n+1) = 2\sum_{p=0}^{n}\left ( \prod_{0\leq r\neq3p\leq 3n}^{{}}\frac{3n+1-r}{3p-r} \right )+ \sum_{p=1}^{n}\left ( \prod_{0\leq r\neq3p-2\leq 3n}^{{}} \frac{3n+1-r}{3p-2-r}\right )$

 after some calculations we get $f(3n+1) =\left ( \binom{3n+1}{0}- \binom{3n+1}{3}+\binom{3n+1}{6}- ... \right )\left ( 2.(-1)^{3n}-1 \right )+1$

 Given $f(3n+1)= 730$ so we have to find $n$ such that $\left ( \binom{3n+1}{0}- \binom{3n+1}{3}+\binom{3n+1}{6}- ... \right )\left ( 2.(-1)^{3n}-1 \right )= 729$

 Lemma: If $p$ is even  $\binom{p}{0}- \binom{p}{3}+ \binom{p}{6}- \cdots = \frac{2^{p+1}sin^{p}\left ( \frac{\pi}{3} \right )(i)^{p}\left ( cos\left ( \frac{p\pi}{3} \right ) \right )}{3}$

 and if  $p$ is odd  $\binom{p}{0}- \binom{p}{3}+ \binom{p}{6}- \cdots = \frac{-2^{p+1}sin^{p}\left ( \frac{\pi}{3} \right )(i)^{p+1}\left ( sin\left ( \frac{p\pi}{3} \right ) \right )}{3}$

 $i$ is $\sqrt{-1}$Using above lemmas we do not get any solution when $n$ is odd, but when $n$ is even $3n+1=13$ satisfies the required condition, hence $n=4$

 


 


\end{enumerate}